% =============================================================================
% SECTION 8 — DISCUSSION. Interprets §7 into a deployer decision rule. Genre-B:
% the "so what" chapter. Must NOT reintroduce claims §7 scoped away (no speedup
% headline; recompile time is measured sub-second R_static~0.3s, framed as
% regeneration/re-audit footprint not wall-clock; frontier is single-run).
% =============================================================================

\section{Discussion}\label{sec:discussion}

The measurements of Section~\ref{sec:evaluation} and the security analysis of
Section~\ref{sec:security} together answer the question this thesis set out to
ask. That question concerns whether a general-purpose zkVM is a viable
substrate for post-quantum anonymous credentials, and at what cost, rather than
how fast the zkVM runs in isolation. We begin with the sharpest of these
findings, the security--deployability tension, then turn the surrounding cost
picture into guidance a deployer can act on.

\subsection{The Security--Deployability Tension}\label{ssec:tension-discussion}

The sharpest finding is the tension of
Section~\ref{sec:security}, which outweighs any single timing figure. On the
target hardware at $\lambda=80$, the proof that is feasible to generate is not
zero-knowledge, and the proof that is zero-knowledge, the one BDEC's anonymity
guarantee requires, is not feasible to generate within the memory budget. A
deployer who needs anonymity on commodity hardware today cannot, on the evidence
here and with the current toolchain versions, have it for this workload. That
statement is scoped to the zkVM stack: the static-circuit regime produces a
zero-knowledge post-quantum proof of the Loquat single-verification building block in
about $22$ minutes (mean of three runs, $21.4$--$22.9$) on the same machine
(\texttt{docs/measurements/audit5\_campaign\_20260612}), at the per-change
regeneration cost this section weighs. More
memory would clear the resource obstruction without restoring the post-quantum
anonymity that motivates the construction, for the substrate-specific reasons
established in Section~\ref{sec:security}.
Recovering it requires a post-quantum-sound zero-knowledge wrap, a STARK- or
lattice-based recursive argument, which the toolchains do not yet provide. We
locate this obstruction along two axes at once, resource and primitive, and it
holds independently of how favourable our timing numbers are and of the open
AIR-soundness Assumption~\ref{asm:air}, which conditions the feasibility figure rather than
the obstruction (Section~\ref{sec:security}).

\subsection{A Substrate-Selection Rule}\label{ssec:decision}

The two regimes do not dominate each other; they trade different costs, and the
right choice depends on a single property of the deployment: how often the
verification predicate changes relative to how often proofs are generated.

\begin{itemize}
  \item \textbf{When the predicate is fixed at deployment time} and proofs are
  generated frequently, the static-circuit regime is preferable: its
  specialised R1CS amortises across many showings, and it pays its update cost
  only once. The zkVM's per-proof overhead (Section~\ref{sec:evaluation}) is not
  justified when there is nothing to update.
  \item \textbf{When the predicate is chosen at presentation time, or the
  signature or security parameters evolve over the deployment's lifetime}, the
  static-circuit regime pays its update-churn cost (Table~\ref{tab:churn})
  \emph{repeatedly}, and the zkVM regime, which absorbs each such change as a
  guest-program edit, becomes the structurally appropriate substrate, at the
  proving cost quantified in Section~\ref{sec:evaluation}. For a
  non-preprocessing static argument such as Aurora this update-churn cost is a
  regeneration-and-re-audit (governance and assurance) burden rather than a
  large wall-clock penalty, which the next paragraph makes precise.
\end{itemize}

\noindent This rule is made precise by the crossover of
\S\ref{sssec:crossover}: the predicate-change rate $\bigl(U/P\bigr)^{*}$ above
which the zkVM is the cheaper substrate over a deployment lifetime is the ratio
of the zkVM's per-proof penalty to the static circuit's per-change recompilation
cost. The recompilation cost is now measured (Section~\ref{sec:evaluation}): for a
transparent non-preprocessing argument such as Aurora it is
$R_{\mathsf{static}}\approx 0.3$~s (R1CS construction plus a microsecond parameter
setup), and it sits outside the multi-minute prover, which is paid per proof.
At that magnitude, and with the zkVM dearer per proof, the wall-clock crossover
favours the static circuit in any realistic churn regime: against Aurora the zkVM
has no wall-clock flexibility advantage. The flexibility claim we can defend is
the \emph{qualitative} regeneration-and-re-audit footprint of
Table~\ref{tab:churn} (zero on the zkVM side, $\geq 1$ on the static side per
change), which is an engineering-and-assurance cost paid by humans, not prover
time. A wall-clock advantage appears only against a \emph{preprocessing}
baseline, as Section~\ref{sec:evaluation} measures for the transparent
preprocessing argument Fractal (Table~\ref{tab:rstatic-fractal}), whose
per-circuit indexer takes $0.86$--$18.5$~s at small synthetic sizes and is
itself OOM-killed at the BDEC circuit size on $24$~GB, and as would hold in
principle for Groth16, Marlin, or circuit-specific PLONK; the non-preprocessing
Aurora that BDEC uses does not offer one. This zero-regeneration property is partitioned rather
than absolute. A change confined to the verification predicate is a guest-program
edit that reuses the existing universal relation, but a change to the hash family
or to the field width still regenerates and re-audits the bespoke
Griffin-Fp192 precompile, so the decision rule applies to predicate churn and
not to substrate-level changes of the signature itself. What the cost model
licenses is limited: it fixes the \emph{form} of the rule and the \emph{sign}
of every term, but not the numeric location of the crossover for PLUM, which
requires a same-scheme static prover time we do not yet have at PLUM's field
width (\S\ref{sssec:crossover}); the Loquat counterpart is measured in the next
paragraph. So the rule gives directional guidance, telling a
deployer which side a high- or low-churn deployment lies on without yet handing
them a numeric threshold to evaluate.

\noindent The contribution of this thesis to that rule is the
\emph{quantification of both sides}: the zkVM's per-proof cost is now measured
(Tables~\ref{tab:plum-verify}--\ref{tab:bdec}), and the structural form of the
static circuit's update cost is enumerated (Table~\ref{tab:churn}). The static circuit's recompile
\emph{time} is now measured ($R_{\mathsf{static}}\approx 0.3$~s, \S\ref{ssec:churn-eval})
and is sub-second because Aurora is non-preprocessing; the remaining term is a
\emph{same-scheme} per-proof prover time on both substrates. For \emph{Loquat} we
now measure it directly: at single-verification scope, $\lambda=80$, with neither
side zero-knowledge, the static Aurora circuit proves in $3.76$~min against
$44$--$64$~min for the zkVM across three runs each
(Table~\ref{tab:loquat-substrate}), so on this clean
single-axis comparison the zkVM is the dearer per-proof substrate by
$12$--$17\times$; the sign of the crossover for Loquat is therefore fixed, with
the static circuit the cheaper prover. For PLUM the sign remains a conjecture: PLUM
over its $199$-bit field has no Aurora harness (the harness is $127$-bit only), so
we cannot yet repeat the measurement on the thesis's headline scheme. We expect the
same direction, and the Loquat measurement is direct evidence for it, but we have
not measured PLUM.

\subsection{What Generalises and What Does Not}\label{ssec:generalise}

The arithmetic-mismatch mechanism (Section~\ref{sec:theoretical}) and the
update-churn dimension (Definition~\ref{def:update-churn}) are
substrate-independent: they apply to any large-field algebraic-hash
post-quantum signature of this form verified in any small-field zkVM, not only
to PLUM in RISC~Zero/SP1. The
specific magnitudes (Section~\ref{sec:evaluation}) do not generalise: they are
tied to PLUM's $199$-bit field, the M5~Pro's $24$~GB envelope, and the CPU-bound
provers of the two zkVMs studied. We drew the boundary deliberately, so that the
shape of the cost carries to other pairings of a large-field signature with a
small-field zkVM while the numbers stay pinned to the configuration measured
here.

\subsection{Comparison to Prior Evaluations}\label{ssec:prior-eval}

As anticipated in Section~\ref{sec:related}, the prior-evaluation record omits
two axes: the cost of a \emph{change} to the verification predicate, and the
consumer-hardware proving frontier for an algebraic-hash post-quantum signature.
We do not restate that gap here; instead we report what our measurements add to
the record. For predicate-change cost, Table~\ref{tab:churn} fixes its structural
form (zero regenerations on the zkVM side against $\geq 1$ per change on the
static side) and \S\ref{ssec:churn-eval} bounds its wall-clock magnitude; for the
frontier, \S\ref{ssec:resource-frontier} locates the binding constraint as the
zero-knowledge wrap rather than the base proof as built and measured (the
base-proof cost is the honest-path cost of the chip and carries
Assumption~\ref{asm:air}). Where our per-proof timings
overlap the existing record, they agree in order of magnitude with the
hitherto-qualitative observation that algebraic-hash workloads are expensive
inside small-field zkVMs, now quantified at the verification level
(Tables~\ref{tab:plum-verify}--\ref{tab:ablation}). The next section draws these
threads together into the thesis's overall conclusion.
